\newcommand{\cert}{\textit{cert}}
\newcommand{\fals}{\textit{fals}}

\Problem{Expressió booleana completament parentitzada}

\Statement

Feu un programa que llegeixi una expressió booleana completament parentitzada,
i que escrigui el resultat d'avaluar-la.
Els operadors admissibles són \emph{conjunció} i \emph{disjunció}
(binaris) i \emph{negació} (unari).


\Input


Els valors booleans es representen per \q{0} i \q{1}.
La conjunció, la disjunció i la negació es representen amb
 \q{*}, \q{+} i \q{!} respectivament.

\medskip

L'entrada és una expressió completament parentitzada.
Això vol dir que sempre apareixen parèntesis
al voltant de les expressions amb operadors binaris.
Per exemple,
\begin{itemize*}

    \item l'expressió $\cert \land \fals$ vindria donada així: \texttt{(1*0)}
    \item l'expressió $\fals \lor (\cert \land \fals)$ vindria donada així: \texttt{(0+(1*0))}
    \item l'expressió $\lnot(\fals \lor (\cert \land \fals))$ vindria donada així: \texttt{!(0+(1*0))}

\end{itemize*}


\Output

Cal escriure el resultat d'avaluar l'expressió en una línia.


\Hint

Fixeu-vos que una expressió és d'una d'aquestes possibilitats:
\begin{itemize*}
\item o bé és directament un dígit zero o u;
\item o bé és l'operador unari de negació, seguit d'una expressió;
\item o bé és un parèntesi obert, seguit d'una expressió, d'un operador binari,
d'una altra expressió, i d'un parèntesi tancat.
\end{itemize*}

\Sample

